\begin{textsample}{POS Dim 4 – human – Score 5.00 – es/es\_ef\_anos\_iniciais\_linguagens\_lp\_3\_p4.txt}  \label{ex:f4_pos_013}
\textbf{Escrita} Compreende a produção de \textbf{textos} não literários verbais, verbo-visuais, multimodais, midiáticos, além de estratégias de planejamento, revisão, reescrita, e avaliação, adequados ao contexto de produção, ao uso da variedade linguística apropriada a esse contexto, aos enunciadores do discurso ( autor e possível leitor ), ao gênero \textbf{textual}, ao suporte e à esfera de circulação.
Abrange ainda a edição de \textbf{textos} nos meios digitais.
A \textbf{escrita} é um processo complexo que exige um projeto de \textbf{texto} organizado, a partir de um gênero e etapas de reescrita, além de distribuição gráfica e marcas de segmentação.
Por isso, este \textbf{eixo} está \textbf{associado} a todos os outros.
Nos dois primeiros anos, temos a importância do professor escriba, que registra a produção individual e coletiva, permitindo que os estudantes tenham a vivência da \textbf{escrita} e a experiência da criação, antes de se apropriarem do sistema de \textbf{escrita} alfabético.

% matched lemmas: associar, eixo, escrita, texto, textual
\end{textsample}
